% !TeX root = ../tesis.tex

% Thesis Abstract -----------------------------------------------------

%\begin{abstractslong}    %uncommenting this line, gives a different abstract heading
\begin{abstracts}        %this creates the heading for the abstract page
\addcontentsline{toc}{chapter}{\protect\numberline{}Resumen}


Los eritrocitos presentan propiedades ópticas que dependen de su geometría, tamaño, composición y orientación, por lo que alteraciones morfológicas y variaciones en la concentración de hemoglobina pueden modificar la manera en que absorben y esparcen la luz. Estas alteraciones ha motivado el uso de técnicas ópticas, como la citometría de flujo, para la caracterización y diferenciación celular. Sin embargo, gran parte de los estudios numéricos del esparcimiento de luz por eritrocitos se han concentrado en luz coherente a una sola longitud de onda fuera de las principales bandas de absorción de la hemoglobina ---Soret y bandas Q---, principal componente de los eritrocitos que domina su respuesta óptica. En este trabajo de tesis de licenciatura se estudian numéricamente las propiedades ópticas de un eritrocito sano, así como de dos de sus variaciones asociadas a enfermedades (el esferocito y el microcito) en una ventana espectral de 400 a 600 nm. Las distintas  morfologías se modelaron como sólidos de revolución a partir de óvalos de Cassini y sus propiedades ópticas se incorporaron a través de funciones dieléctricas dependientes de la concentración de hemoglobina. Las funciones dieléctricas en el rango espectral de interés se obtuvieron a partir de las relaciones de Kramers-Kronig sustractivas a datos experimentales en rangos espectrales más reducidos. Las simulaciones numéricas se realizaron mediante el método \textit{Finite Integration Technique} implementadas en CST Studio Suite. Se determinaron los parámetros óptimos para garantizar la convergencia numérica al contrastar con la solución analítica de teoría de con Mie, obteniendo errores menores al 5\%. Además, se reprodujeron resultados reportados en la literatura, mostrando fiabilidad en los parámetros seleccionados para la simulación. Como resultados, se observa que la banda de Soret, alrededor de 431 nm,  presenta las mayores diferencias en la eficiencia de extinción entre las morfologías consideradas. Aunque el esparcimiento se concentra predominantemente en la dirección frontal, las mayores diferencias relativas aparecen en regiones laterales y posteriores.  Finalmente, el análisis mediante ventanas angulares permitió identificar una región del ángulo de esparcimiento alrededor de $130^\circ$ en la que el contraste entre eritrocitos sanos y patológicos aumenta de manera consistente para las tres longitudes de onda analizadas. A partir de los resultados obtenidos se identifican condiciones espectrales y angulares potencialmente útiles para futuros esquemas de diferenciación celular mediante citometría de flujo.


  

\end{abstracts}
%\end{abstractlongs}


% ----------------------------------------------------------------------